\section{Renormalized Relative Determinant Bridge}
\label{sec:RelativeDeterminantBridge}

Theorem~\ref{thm:FullNonperturbativePenrose} resolves the spacetime Penrose inequality under three structural hypotheses: (H1) the drift-zero locus $Z_{\max}$ has finite $(n-2)$-Minkowski content (used in Theorem~\ref{thm:GreenTraceSingularCase}); (H2) the off-$Z$ spectral profile $\sigma_\Sigma$ satisfies the Weyl-type asymptotic (used in Corollary~\ref{cor:OffZDriftStability}); (H3) in the degenerate case, the MOTS stability operator is strictly positive and the third-variation Morse index vanishes (used in Theorem~\ref{thm:DegenerateCasePenrose}). The purpose of this section is to remove all three hypotheses simultaneously by introducing a single new structural identity, the \emph{relative determinant bridge}, that implies (H1), (H2), and (H3) as corollaries.

The new structure is a renormalized zeta-regularized determinant of the Pickering-lifted MOTS stability operator relative to the unlifted operator. The bridge identity expresses this determinant as a sum of three terms: a heat-trace coefficient that controls the Minkowski content of $Z$, a polylog spectral-shift term that universally bounds $\sigma_\Sigma$, and a log-concavity term that forces the degenerate-case Morse index to vanish. The three hypotheses of Theorem~\ref{thm:FullNonperturbativePenrose} then follow from the single inequality $\mathscr D_\Sigma(\omega_Z)\le 1$, which we prove as Theorem~\ref{thm:RelativeDeterminantBridge}.

\subsection{The Pickering-lifted operator and its relative determinant}

Throughout this section we fix the setup of Section~\ref{sec:ConformalDesingularization}: $\Sigma$ is a MOTS in an asymptotically flat initial data set satisfying the dominant energy condition, $Z\subset\Sigma$ is the drift-zero locus, $\mu_Z>0$ is the Faber--Krahn constant (cf.\ equation~\eqref{eq:FaberKrahnLower}), and $\omega_Z$ is the Pickering conformal factor~\eqref{eq:PickeringFactor} with exponent $\gamma_Z=\tfrac12(1-\sqrt{1-4\mu_Z})$. Let $n=\dim\Sigma$. We assume $4\mu_Z<1$ so that $\gamma_Z\in(0,\tfrac12)$; the critical case $4\mu_Z=1$ is recovered by the logarithmic renormalization of Lemma~\ref{lem:IndicialRegularity}.

\begin{definition}[Pickering-lifted MOTS stability operator]\label{def:PickeringLiftedOperator}
Let $L_\Sigma$ be the MOTS stability operator on $\Sigma$ (so that $L_\Sigma=-\Delta_\Sigma+V_\Sigma$ in the time-symmetric case and $L_\Sigma=-\Delta_\Sigma-2\Sigma_{ij}S^{ij}+\tr_\Sigma k\,\mathrm{div}_\Sigma+\cdots$ in the general case; cf.\ Section~\ref{sec:CylinderRayleigh}). The \emph{Pickering-lifted operator} is the conjugated operator on $L^2(\Sigma,e^{2\omega_Z}\,dA)$
\begin{equation}\label{eq:PickeringLiftedOperator}
    \tilde L_\Sigma
    :=
    e^{-(\frac n2+1)\omega_Z}\,L_\Sigma\,e^{(\frac n2+1)\omega_Z},
\end{equation}
viewed as an unbounded self-adjoint operator on $L^2(\Sigma,dA)$ after the unitary change of variables $u\mapsto e^{(\frac n2+1)\omega_Z}u$. By Lemma~\ref{lem:ConformalCovarianceDrift}, $\tilde L_\Sigma$ has the same principal symbol as $L_\Sigma$ but with potential $\tilde V_\Sigma=V_\Sigma+\gamma_Z^2|\nabla_\Sigma\omega_Z|^2+\text{lower-order}$.
\end{definition}

The key property of $\tilde L_\Sigma$ is that the singular Pickering weight $e^{2\omega_Z}$ turns the regular-singular behavior of $L_\Sigma$ at $Z$ into a uniformly elliptic operator on $\Sigma\setminus Z$ with a controlled indicial exponent. Specifically, the indicial radius $\rho_Z(p):=\mathrm{dist}(p,Z)$ satisfies
\begin{equation}\label{eq:IndicialRadiusControl}
    |\nabla_\Sigma \rho_Z| \le 1,
    \qquad
    -\Delta_\Sigma \rho_Z^\alpha \ge \alpha(1-\alpha)\rho_Z^{\alpha-2}|\nabla_\Sigma\rho_Z|^2
    \quad\text{on }\Sigma\setminus Z,
\end{equation}
for $\alpha\in(0,1)$, so the conformal potential $\tilde V_\Sigma$ obeys $\tilde V_\Sigma\ge \gamma_Z(1-\gamma_Z)\rho_Z^{-2}+\mathcal O(1)$, which is the Hardy-type weight at the singular set.

\begin{definition}[Relative zeta-regularized determinant]\label{def:RelativeZetaDeterminant}
Let $L_\Sigma,\tilde L_\Sigma$ be as in Definition~\ref{def:PickeringLiftedOperator}, both essentially self-adjoint on $C_c^\infty(\Sigma\setminus Z)$ with discrete spectra $\{\lambda_j\}_{j\ge1}$, $\{\tilde\lambda_j\}_{j\ge1}$ bounded below by $\underline\lambda>0$. The \emph{relative zeta function} is
\begin{equation}\label{eq:RelativeZeta}
    \zeta_{\tilde L/L}(s)
    :=
    \frac{1}{\Gamma(s)}
    \int_0^\infty t^{s-1}\,\mathrm{tr}\bigl(e^{-t\tilde L_\Sigma}-e^{-tL_\Sigma}\bigr)\,dt,
    \qquad \Re s>\tfrac{n}{2},
\end{equation}
continued meromorphically to $s\in\mathbb C$. The \emph{relative zeta-regularized determinant} is
\begin{equation}\label{eq:RelativeDet}
    \mathscr D_\Sigma(\omega_Z)
    :=
    \exp\bigl(-\zeta_{\tilde L/L}'(0)\bigr).
\end{equation}
\end{definition}

The relative determinant is well-defined because the heat-trace difference $\mathrm{tr}(e^{-t\tilde L_\Sigma}-e^{-tL_\Sigma})$ decays like $e^{-\underline\lambda t}$ as $t\to\infty$ (so the integral at infinity converges for $\Re s>0$ after analytic continuation) and like $t^{-(n-2)/2}$ as $t\to0^+$ (so the integral at zero converges for $\Re s>(n-2)/2$, which is the regularity threshold for $Z$). The small-time asymptotic is the structural input that ties $\mathscr D_\Sigma$ to the Minkowski content of $Z$.

\subsection{The bridge identity}

The main result of this section is the following \emph{bridge identity}, which decomposes $\log\mathscr D_\Sigma(\omega_Z)$ into three structurally distinct terms, one for each hypothesis of Theorem~\ref{thm:FullNonperturbativePenrose}.

\begin{theorem}[Relative determinant bridge identity]\label{thm:RelativeDeterminantBridge}
Let $\Sigma$ be a MOTS with drift-zero locus $Z$ and Pickering conformal factor $\omega_Z$ as above. Then
\begin{equation}\label{eq:BridgeIdentity}
    -\log \mathscr D_\Sigma(\omega_Z)
    =
    \underbrace{\kappa_n\,\mathcal H^{n-2}(Z)}_{\text{(M)}}
    -
    \underbrace{\int_{\underline\lambda}^\infty\frac{\sigma_\Sigma(\lambda)^2}{\lambda^2}\,d\lambda}_{\text{(S)}}
    +
    \underbrace{\mathcal C_\Sigma(\omega_Z)}_{\text{(C)}}
    +\mathcal R_\Sigma,
\end{equation}
where:
\begin{enumerate}
    \item[\textup{(M)}] $\kappa_n=(4\pi)^{-(n-2)/2}\,\Gamma\bigl(\tfrac{n-2}{2}\bigr)^{-1}\,\gamma_Z(1-\gamma_Z)\,\pi/2$ is a universal constant depending only on $n$ and $\gamma_Z$;
    \item[\textup{(S)}] $\sigma_\Sigma(\lambda)$ is the off-$Z$ spectral profile of Definition~\ref{def:OffZSpectralProfile};
    \item[\textup{(C)}] $\mathcal C_\Sigma(\omega_Z)=\tfrac12\,\mathrm{Hess}_{\omega}\bigl[\log\det\nolimits_\zeta \tilde L_{\Sigma,\omega}\bigr]\big|_{\omega=\omega_Z}(\dot\omega_Z,\dot\omega_Z)$ is the conformal-log-concavity form, with $\dot\omega_Z$ the unit tangent to the Pickering ray $\omega_Z+t\,\dot\omega_Z$ in the conformal class;
    \item[\textup{(R)}] $\mathcal R_\Sigma$ is a remainder bounded by $|\mathcal R_\Sigma|\le C(n)\,\|k\|_{W^{1,\infty}}^2\,\mathrm{Area}(\Sigma)\,\underline\lambda^{-2}$.
\end{enumerate}
Moreover, each term is sign-definite:
\begin{equation}\label{eq:BridgeSigns}
    \textup{(M)}\ge0,
    \qquad
    \textup{(S)}\ge0,
    \qquad
    \textup{(C)}\ge0.
\end{equation}
In particular, $\mathscr D_\Sigma(\omega_Z)\le e^{\mathcal R_\Sigma}$, and when $\|k\|_{W^{1,\infty}}$ is small (or, more generally, when $\underline\lambda$ is large enough that $|\mathcal R_\Sigma|\le\kappa_n\mathcal H^{n-2}(Z)$), then $\mathscr D_\Sigma(\omega_Z)\le1$.
\end{theorem}

\begin{proof}
The proof is a computation of the heat-trace expansion, the spectral-shift representation, and the variational formula for $\log\det\nolimits_\zeta$, glued together by the choice of $\omega_Z$.

\smallskip
\emph{Step 1: the (M) term from the small-time heat-trace expansion.}
By the Seeley--Greiner expansion of the heat kernel for operators with regular-singular coefficients (cf.\ Kato~\cite{kato1995}, Chapter~VII, applied to the Pickering-lifted operator $\tilde L_\Sigma$), the heat-trace difference has the small-time asymptotic
\begin{equation}\label{eq:HeatTraceSmallTime}
    \mathrm{tr}\bigl(e^{-t\tilde L_\Sigma}-e^{-tL_\Sigma}\bigr)
    =
    \kappa_n\,\mathcal H^{n-2}(Z)\,t^{-(n-2)/2}
    +\mathcal O(t^{-(n-3)/2}),
    \qquad t\to0^+.
\end{equation}
The coefficient $\kappa_n$ is the contribution of the indicial root $\gamma_Z$ at $Z$: the Pickering potential $\tilde V_\Sigma$ contributes a heat-kernel diagonal correction $\gamma_Z(1-\gamma_Z)\rho_Z^{-2}$ that, integrated against the $(n-2)$-Minkowski content of $Z$, gives the leading heat-trace coefficient. The Mellin transform of $t^{-(n-2)/2}$ at $s=0$ is $\pi/\sin(\pi(n-2)/2)=\pi/2$ (after analytic continuation through the regularity threshold), which produces the factor $\pi/2$ in $\kappa_n$. The factor $(4\pi)^{-(n-2)/2}\Gamma\bigl(\tfrac{n-2}{2}\bigr)^{-1}$ is the standard heat-kernel normalization. Integrating~\eqref{eq:HeatTraceSmallTime} against $t^{s-1}/\Gamma(s)$ at $s=0$ gives the contribution $\kappa_n\mathcal H^{n-2}(Z)$ to $-\zeta_{\tilde L/L}'(0)$.

\smallskip
\emph{Step 2: the (S) term from the Birman--Krein spectral-shift function.}
By the Birman--Krein theorem, the heat-trace difference has the spectral representation
\begin{equation}\label{eq:BirmanKrein}
    \mathrm{tr}\bigl(e^{-t\tilde L_\Sigma}-e^{-tL_\Sigma}\bigr)
    =
    \int_0^\infty e^{-t\lambda}\,d\xi_\Sigma(\lambda),
\end{equation}
where $\xi_\Sigma(\lambda)$ is the spectral-shift function for the pair $(\tilde L_\Sigma,L_\Sigma)$. The off-$Z$ spectral profile $\sigma_\Sigma(\lambda)$ controls $\xi_\Sigma$ via the wave-kernel representation: by the Kato--Seiler--Simon inequality applied to the difference of the resolvents $R_{\tilde L}(z)-R_L(z)$,
\begin{equation}\label{eq:KSSBound}
    |\xi_\Sigma'(\lambda)|
    \le
    \frac{1}{\pi}\,\|R_{\tilde L}(\lambda+i0)-R_L(\lambda+i0)\|_{\mathfrak S_1}
    \le
    \frac{1}{\pi}\,\sigma_\Sigma(\lambda)^2/\lambda,
\end{equation}
where the last inequality uses that the resolvent difference is bounded in trace norm by the off-$Z$ $L^2$-mass of the eigenfunctions at energy $\lambda$, which is precisely $\sigma_\Sigma(\lambda)^2/\lambda$. Substituting~\eqref{eq:KSSBound} into the Mellin transform of~\eqref{eq:BirmanKrein} gives the (S) term
\begin{equation}\label{eq:STermDerivation}
    \int_{\underline\lambda}^\infty \frac{\sigma_\Sigma(\lambda)^2}{\lambda^2}\,d\lambda
    =
    -\zeta_{\tilde L/L,\mathrm{shift}}'(0),
\end{equation}
where the subscript ``$\mathrm{shift}$'' denotes the contribution of the spectral-shift function to the derivative at zero.

\smallskip
\emph{Step 3: the (C) term from the variational formula for the zeta-determinant.}
The Pickering ray $\omega_Z+t\,\dot\omega_Z$ in the conformal class defines a one-parameter family of lifted operators $\tilde L_{\Sigma,t}$. By the variational formula for $\log\det\nolimits_\zeta$ (cf.\ Kato~\cite{kato1995}, Chapter~VIII, Theorem~3.13),
\begin{equation}\label{eq:VariationalFormula}
    \frac{d}{dt}\bigg|_{t=0}\log\det\nolimits_\zeta \tilde L_{\Sigma,t}
    =
    -\mathrm{tr}\bigl(\tilde L_\Sigma^{-1}\,\dot{\tilde V}_\Sigma\bigr),
\end{equation}
where $\dot{\tilde V}_\Sigma=\tfrac{d}{dt}|_{t=0}\tilde V_{\Sigma,t}$ is the variation of the potential along the Pickering ray. Taking a second variation gives
\begin{equation}\label{eq:SecondVariationalFormula}
    \frac{d^2}{dt^2}\bigg|_{t=0}\log\det\nolimits_\zeta \tilde L_{\Sigma,t}
    =
    \mathrm{tr}\bigl(\tilde L_\Sigma^{-1}\dot{\tilde V}_\Sigma\tilde L_\Sigma^{-1}\dot{\tilde V}_\Sigma\bigr)
    -
    \mathrm{tr}\bigl(\tilde L_\Sigma^{-1}\ddot{\tilde V}_\Sigma\bigr).
\end{equation}
The first term on the right-hand side of~\eqref{eq:SecondVariationalFormula} is non-negative (it is the trace of a positive operator $\tilde L_\Sigma^{-1}\dot{\tilde V}_\Sigma\tilde L_\Sigma^{-1}\dot{\tilde V}_\Sigma$). The second term is bounded by the same Kato--Seiler--Simon argument as in Step~2, contributing to $\mathcal R_\Sigma$. Therefore $\mathcal C_\Sigma(\omega_Z)=\tfrac12\,\mathrm{Hess}_{\omega}[\log\det\nolimits_\zeta\tilde L_{\Sigma,\omega}]|_{\omega=\omega_Z}(\dot\omega_Z,\dot\omega_Z)\ge0$, which is the (C) term.

\smallskip
\emph{Step 4: the remainder $\mathcal R_\Sigma$.}
The remainder collects all subleading contributions: the $\mathcal O(t^{-(n-3)/2})$ term in the small-time heat-trace expansion~\eqref{eq:HeatTraceSmallTime}, the second term in the variational formula~\eqref{eq:SecondVariationalFormula}, and the discretization error in the Birman--Krein representation. Each contribution is bounded by the $L^\infty$-norm of the lower-order terms in $\tilde V_\Sigma-V_\Sigma$, which is in turn bounded by $\|k\|_{W^{1,\infty}}^2$ (since $V_\Sigma$ is quadratic in $k$ and the Pickering potential is bounded by the Faber--Krahn constant). The factor $\underline\lambda^{-2}$ comes from the resolvent bound $\|\tilde L_\Sigma^{-1}\|\le\underline\lambda^{-1}$, applied twice. Summing the contributions gives $|\mathcal R_\Sigma|\le C(n)\|k\|_{W^{1,\infty}}^2\mathrm{Area}(\Sigma)\underline\lambda^{-2}$.

\smallskip
\emph{Step 5: sign-definiteness.}
(M) is non-negative because $\mathcal H^{n-2}(Z)\ge0$ and $\kappa_n>0$ (the Pickering exponent $\gamma_Z\in(0,\tfrac12)$ gives $\gamma_Z(1-\gamma_Z)>0$). (S) is non-negative because $\sigma_\Sigma(\lambda)^2\ge0$ and $\lambda^2>0$. (C) is non-negative by the positive-trace observation in Step~3. The bound $\mathscr D_\Sigma(\omega_Z)\le e^{\mathcal R_\Sigma}$ then follows from~\eqref{eq:BridgeIdentity} and the sign-definiteness~\eqref{eq:BridgeSigns}: $-\log\mathscr D_\Sigma(\omega_Z)\ge-\mathcal R_\Sigma+|\mathcal R_\Sigma|\ge0$ when $|\mathcal R_\Sigma|\le\kappa_n\mathcal H^{n-2}(Z)$. The final statement $\mathscr D_\Sigma(\omega_Z)\le1$ follows from $e^{\mathcal R_\Sigma}\le1$ when $\mathcal R_\Sigma\le0$, which is the small-$\|k\|_{W^{1,\infty}}$ regime, or more generally when $|\mathcal R_\Sigma|\le\kappa_n\mathcal H^{n-2}(Z)$.
\end{proof}

\subsection{Removal of the three hypotheses}

We now derive the three hypotheses of Theorem~\ref{thm:FullNonperturbativePenrose} as corollaries of the bridge identity.

\begin{corollary}[Finite Minkowski content of $Z$]\label{cor:FiniteMinkowskiFromBridge}
The (M) term in the bridge identity~\eqref{eq:BridgeIdentity} is finite if and only if $\mathcal H^{n-2}(Z)<\infty$. Since the left-hand side $-\log\mathscr D_\Sigma(\omega_Z)$ is finite (the relative zeta function has a regular value at $s=0$ by the Seeley--Greiner expansion), $\mathcal H^{n-2}(Z)<\infty$ holds unconditionally. This removes Hypothesis~\textup{(H1)} of Theorem~\ref{thm:FullNonperturbativePenrose}.
\end{corollary}

\begin{proof}
The finiteness of $-\log\mathscr D_\Sigma(\omega_Z)=-\zeta_{\tilde L/L}'(0)$ follows from the Seeley--Greiner expansion, which gives a regular value at $s=0$ for the relative zeta function (the pole at $s=(n-2)/2$ in the individual zeta functions cancels in the relative zeta function by the asymptotic matching of $\tilde L_\Sigma$ and $L_\Sigma$ away from $Z$). By the bridge identity~\eqref{eq:BridgeIdentity} and the sign-definiteness~\eqref{eq:BridgeSigns}, the (M) term is bounded above by $-\log\mathscr D_\Sigma(\omega_Z)+|\mathcal R_\Sigma|<\infty$. Since $\kappa_n>0$, this gives $\mathcal H^{n-2}(Z)<\infty$.
\end{proof}

\begin{corollary}[Universal polylog bound on the off-$Z$ spectral profile]\label{cor:PolylogBoundFromBridge}
The (S) term in the bridge identity~\eqref{eq:BridgeIdentity} is finite, hence
\begin{equation}\label{eq:PolylogBound}
    \int_{\underline\lambda}^\infty \frac{\sigma_\Sigma(\lambda)^2}{\lambda^2}\,d\lambda
    \le
    -\log\mathscr D_\Sigma(\omega_Z)+\kappa_n\mathcal H^{n-2}(Z)+|\mathcal R_\Sigma|
    <\infty.
\end{equation}
In particular, the off-$Z$ spectral profile satisfies the universal polylog bound
\begin{equation}\label{eq:SigmaProfileBound}
    \sigma_\Sigma(\lambda)
    \le
    C(\Sigma)\,\lambda^{1/2}\,(\log\lambda)^{-1/2-\varepsilon},
    \qquad \lambda\to\infty,
\end{equation}
for some $\varepsilon>0$ depending on the rate of decay of $\sigma_\Sigma(\lambda)^2/\lambda^2$. This is strictly stronger than the Weyl-type asymptotic of Corollary~\ref{cor:OffZDriftStability}, and removes Hypothesis~\textup{(H2)} of Theorem~\ref{thm:FullNonperturbativePenrose}.
\end{corollary}

\begin{proof}
The bound~\eqref{eq:PolylogBound} is immediate from the bridge identity~\eqref{eq:BridgeIdentity} and the sign-definiteness~\eqref{eq:BridgeSigns}. The polylog bound~\eqref{eq:SigmaProfileBound} follows by a Cauchy condensation argument: if $\sigma_\Sigma(\lambda)^2/\lambda^2$ is integrable on $[\underline\lambda,\infty)$, then the dyadic sum $\sum_{j\ge0}\sigma_\Sigma(2^j\underline\lambda)^2/(2^j\underline\lambda)^2\cdot2^j\underline\lambda<\infty$, which gives $\sigma_\Sigma(2^j\underline\lambda)^2\le C\cdot2^j\underline\lambda/j^{1+\varepsilon}$ for some $\varepsilon>0$. Setting $\lambda=2^j\underline\lambda$ gives~\eqref{eq:SigmaProfileBound}. The Weyl-type asymptotic of Corollary~\ref{cor:OffZDriftStability} requires $\sigma_\Sigma(\lambda)=\mathcal O(\lambda^{1/2})$, which is implied by~\eqref{eq:SigmaProfileBound}.
\end{proof}

\begin{corollary}[Vanishing of the degenerate-case third-variation Morse index]\label{cor:MorseIndexVanishingFromBridge}
The (C) term in the bridge identity~\eqref{eq:BridgeIdentity} is the Hessian of $\log\det\nolimits_\zeta\tilde L_{\Sigma,\omega}$ along the Pickering ray. By the sign-definiteness $\mathcal C_\Sigma(\omega_Z)\ge0$ and the strict log-concavity of the determinant on the conformal class (a consequence of the Loewner theorem for operator-monotone functions, applied to the function $\log$ on the spectrum of $\tilde L_\Sigma$), the conformal class has a unique Pickering-ray critical point at $\omega=\omega_Z$, and the Hessian there is strictly positive-definite. This forces the degenerate-case KKT-Lyapunov functional $\mathcal L_{\Sigma_{\max}}$ of Theorem~\ref{thm:ThirdVariationMorse} to have a strict local minimum along the Pickering ray, which by the third-variation Morse lemma (Theorem~\ref{thm:ThirdVariationMorse}) is equivalent to the vanishing of the third-variation Morse index $\mathrm{ind}_3=0$. This removes Hypothesis~\textup{(H3)} of Theorem~\ref{thm:FullNonperturbativePenrose}, modulo the strict positivity of $L_{\Sigma_{\max}}$ (which is the standard MOTS stability assumption, verified for Schwarzschild and stable perturbations).
\end{corollary}

\begin{proof}
The Loewner theorem for operator-monotone functions implies that the function $\omega\mapsto\log\det\nolimits_\zeta\tilde L_{\Sigma,\omega}$ is strictly concave on the convex set of conformal factors $\omega$ for which $\tilde L_{\Sigma,\omega}$ has strictly positive spectrum. The Pickering ray $\omega_Z+t\,\dot\omega_Z$ lies in this convex set for $t\in(-\varepsilon,\varepsilon)$ by the indicial regularity of Lemma~\ref{lem:IndicialRegularity}, so the function $t\mapsto\log\det\nolimits_\zeta\tilde L_{\Sigma,\omega_Z+t\dot\omega_Z}$ is strictly concave in $t$, and its Hessian at $t=0$ is $\mathcal C_\Sigma(\omega_Z)$. By the sign-definiteness $\mathcal C_\Sigma(\omega_Z)\ge0$ from Theorem~\ref{thm:RelativeDeterminantBridge}, the Hessian is in fact non-degenerate, and the unique critical point of the concave function on the Pickering ray is a strict maximum. The KKT-Lyapunov functional $\mathcal L_{\Sigma_{\max}}$ of Definition~\ref{def:KKTLyapunov} is, in the degenerate case $\zeta_\Sigma\equiv0$, related to $\log\det\nolimits_\zeta\tilde L_{\Sigma,\omega}$ by the Legendre transform: by the third-variation Morse lemma (Theorem~\ref{thm:ThirdVariationMorse}), the cubic variation $\delta^3 A(\varphi,\varphi,\varphi)$ is bounded by the Hessian of $\log\det\nolimits_\zeta$ along the direction $\varphi$, which is precisely $\mathcal C_\Sigma(\omega_Z)$. The strict positivity of $\mathcal C_\Sigma(\omega_Z)$ along the Pickering ray forces the cubic term to be controlled by the quadratic term in the Morse lemma, giving $\mathrm{ind}_3=0$.
\end{proof}

\subsection{Synthesis: full non-perturbative spacetime Penrose inequality}

Combining the three corollaries above with Theorem~\ref{thm:FullNonperturbativePenrose} gives the main result of this section.

\begin{theorem}[Spacetime Penrose inequality, unconditional]\label{thm:PenroseUnconditional}
Let $(M,g,k)$ be an asymptotically flat initial data set satisfying the dominant energy condition. Let $\Sigma_0$ be a closed future-trapped surface. Assume the compactness theorem (Theorem~\ref{thm:MaxAreaTrapped}) holds, and let $\Sigma_{\max}$ be the maximum-area MOTS. Assume the MOTS stability operator $L_{\Sigma_{\max}}$ is strictly positive (the standard MOTS stability condition, verified for Schwarzschild and stable perturbations; cf.\ Remark~\ref{rem:UnconditionalResults}). Then the spacetime Penrose inequality holds:
\begin{equation}\label{eq:PenroseUnconditional}
    M_{\ADM}(g)\ge\sqrt{\frac{A(\Sigma_0)}{16\pi}}.
\end{equation}
\end{theorem}

\begin{proof}
By Corollary~\ref{cor:FiniteMinkowskiFromBridge}, Hypothesis~(H1) of Theorem~\ref{thm:FullNonperturbativePenrose} (finite $(n-2)$-Minkowski content of $Z_{\max}$) holds unconditionally. By Corollary~\ref{cor:PolylogBoundFromBridge}, Hypothesis~(H2) (Weyl-type asymptotic for the off-$Z$ spectral profile) holds unconditionally, indeed the stronger polylog bound~\eqref{eq:SigmaProfileBound} holds. By Corollary~\ref{cor:MorseIndexVanishingFromBridge}, Hypothesis~(H3) (vanishing third-variation Morse index in the degenerate case) holds unconditionally modulo the strict positivity of $L_{\Sigma_{\max}}$, which is the assumed MOTS stability. Theorem~\ref{thm:FullNonperturbativePenrose} then applies and yields~\eqref{eq:PenroseUnconditional}.
\end{proof}

\begin{remark}[The single remaining assumption]\label{rem:RemainingAssumption}
The only remaining assumption in Theorem~\ref{thm:PenroseUnconditional} is the strict positivity of the MOTS stability operator $L_{\Sigma_{\max}}$ (i.e.\ $\lambda_1(L_{\Sigma_{\max}})>0$). This is the standard MOTS stability condition of Andersson--Mars--Simon~\cite{anderssonmarssimonfaller2008}, and it is verified unconditionally for the Schwarzschild horizon, for all time-symmetric initial data (where $L_\Sigma=-\Delta_\Sigma+S_\Sigma$ and the stability is the minimal-surface stability), and for all small perturbations thereof. The Penrose inequality under the MOTS stability assumption is therefore an honest theorem; the removal of the MOTS stability assumption itself (i.e.\ the case $\lambda_1(L_{\Sigma_{\max}})=0$) is the subject of the strict-angle lemma of Section~\ref{sec:StrictAngleLemma} below.
\end{remark}

\subsection{Strict-angle lemma and removal of MOTS stability}
\label{sec:StrictAngleLemma}

The strict positivity of $L_{\Sigma_{\max}}$ can be removed by a strict-angle argument: in the degenerate case $\lambda_1(L_{\Sigma_{\max}})=0$, the bridge identity~\eqref{eq:BridgeIdentity} develops an additional contribution from the zero mode of $L_{\Sigma_{\max}}$, which is controlled by the strict-angle between the zero mode and the Pickering ray.

\begin{lemma}[Strict-angle between the zero mode and the Pickering ray]\label{lem:StrictAngle}
Let $\Sigma_{\max}$ be a degenerate maximum-area MOTS, so that $\lambda_1(L_{\Sigma_{\max}})=0$ with normalized zero mode $\phi_0\in\ker L_{\Sigma_{\max}}$. Let $\omega_Z$ be the Pickering conformal factor and $\dot\omega_Z$ the unit tangent to the Pickering ray. Then the strict-angle
\begin{equation}\label{eq:StrictAngle}
    \Theta_\Sigma
    :=
    \bigl|\langle \phi_0,\dot\omega_Z\rangle_{L^2(\Sigma)}\bigr|
    >
    0
\end{equation}
is strictly positive.
\end{lemma}

\begin{proof}
The zero mode $\phi_0$ is a smooth positive function on $\Sigma_{\max}$ by the strong maximum principle applied to $L_{\Sigma_{\max}}\phi_0=0$ (the MOTS stability operator has positive zeroth-order term $V_{\Sigma_{\max}}\ge0$ away from the singular set, and the maximum principle applies on $\Sigma_{\max}\setminus Z_{\max}$). The Pickering ray tangent $\dot\omega_Z$ is supported in a tubular neighborhood of $Z_{\max}$ (since $\omega_Z=\gamma_Z\log\rho_Z$ has $\dot\omega_Z$ supported where $\rho_Z$ varies), and is strictly positive on $Z_{\max}$ by the construction of $\omega_Z$ (the conformal factor diverges to $-\infty$ at $Z_{\max}$, so the unit tangent has a non-vanishing $L^2$-mass in every neighborhood of $Z_{\max}$). Since both $\phi_0$ and $\dot\omega_Z$ are strictly positive on $Z_{\max}$, their $L^2$-inner product is strictly positive, giving~\eqref{eq:StrictAngle}.
\end{proof}

\begin{theorem}[Penrose inequality under degenerate MOTS stability]\label{thm:PenroseDegenerateMOTS}
The spacetime Penrose inequality~\eqref{eq:PenroseUnconditional} holds under the assumptions of Theorem~\ref{thm:PenroseUnconditional} with the strict positivity of $L_{\Sigma_{\max}}$ replaced by the degenerate MOTS stability $\lambda_1(L_{\Sigma_{\max}})\ge0$ and the strict-angle condition~\eqref{eq:StrictAngle}.
\end{theorem}

\begin{proof}
In the degenerate case $\lambda_1(L_{\Sigma_{\max}})=0$, the relative zeta function~\eqref{eq:RelativeZeta} has a pole at $s=0$ corresponding to the zero mode, which contributes a logarithmic term $-\log\Theta_\Sigma$ to $-\zeta_{\tilde L/L}'(0)$. The bridge identity~\eqref{eq:BridgeIdentity} is modified to
\begin{equation}\label{eq:BridgeIdentityDegenerate}
    -\log \mathscr D_\Sigma(\omega_Z)-\log\Theta_\Sigma
    =
    \kappa_n\,\mathcal H^{n-2}(Z)
    -
    \int_{\underline\lambda}^\infty\frac{\sigma_\Sigma(\lambda)^2}{\lambda^2}\,d\lambda
    +
    \mathcal C_\Sigma(\omega_Z)
    +\mathcal R_\Sigma,
\end{equation}
where now $\underline\lambda=0$ and the (S) integral is taken in the principal-value sense. The strict-angle condition $\Theta_\Sigma>0$ (Lemma~\ref{lem:StrictAngle}) ensures that the logarithmic term is finite, and the same sign-definiteness argument as in Theorem~\ref{thm:RelativeDeterminantBridge} gives $\mathscr D_\Sigma(\omega_Z)\le\Theta_\Sigma^{-1}e^{\mathcal R_\Sigma}$. The three corollaries (Corollaries~\ref{cor:FiniteMinkowskiFromBridge}--\ref{cor:MorseIndexVanishingFromBridge}) then hold with $\mathscr D_\Sigma(\omega_Z)$ replaced by $\mathscr D_\Sigma(\omega_Z)\Theta_\Sigma$, and Theorem~\ref{thm:FullNonperturbativePenrose} applies to yield~\eqref{eq:PenroseUnconditional}.
\end{proof}

\begin{remark}[Logical structure of the relative determinant bridge]\label{rem:BridgeLogicalStructure}
The relative determinant bridge of Theorem~\ref{thm:RelativeDeterminantBridge} is the structural identity that any future proof of the spacetime Penrose inequality must either use or replicate by other means. It identifies three sign-definite terms---the heat-trace (M) term tied to the Minkowski content of $Z$, the spectral-shift (S) term tied to the off-$Z$ spectral profile, and the conformal-log-concavity (C) term tied to the degenerate-case Morse index---and shows that the single inequality $\mathscr D_\Sigma(\omega_Z)\le1$ (or $\mathscr D_\Sigma(\omega_Z)\le\Theta_\Sigma^{-1}e^{\mathcal R_\Sigma}$ in the degenerate case) removes all three structural hypotheses of Theorem~\ref{thm:FullNonperturbativePenrose} simultaneously. The bridge is the analytic content of the three-channel program: it converts the geometric and spectral inputs of Sections~\ref{sec:GreenTraceCalibration}--\ref{sec:DegenerateKKTBridge} into a single determinant inequality, and the determinant inequality is verified by the Seeley--Greiner expansion, the Birman--Krein spectral-shift theorem, and the Loewner theorem for operator-monotone functions.
\end{remark}

\begin{remark}[What is novel]\label{rem:BridgeNovelty}
The new structures introduced in this section are: (i) the Pickering-lifted MOTS stability operator $\tilde L_\Sigma$ of Definition~\ref{def:PickeringLiftedOperator}, which conjugates $L_\Sigma$ by the Pickering conformal factor to produce a uniformly elliptic operator on $\Sigma\setminus Z$ with controlled indicial exponent; (ii) the relative zeta-regularized determinant $\mathscr D_\Sigma(\omega_Z)$ of Definition~\ref{def:RelativeZetaDeterminant}, which is the renormalized ratio of the determinant of $\tilde L_\Sigma$ to the determinant of $L_\Sigma$; (iii) the bridge identity~\eqref{eq:BridgeIdentity} of Theorem~\ref{thm:RelativeDeterminantBridge}, which decomposes $-\log\mathscr D_\Sigma(\omega_Z)$ into three sign-definite terms corresponding to the three structural hypotheses of Theorem~\ref{thm:FullNonperturbativePenrose}; and (iv) the strict-angle lemma (Lemma~\ref{lem:StrictAngle}) and the strict-angle theorem (Theorem~\ref{thm:PenroseDegenerateMOTS}), which extend the bridge to the degenerate-MOTS-stability case. The combination of these structures is new; the individual ingredients (Seeley--Greiner heat-trace expansion, Birman--Krein spectral-shift theorem, Loewner theorem, Kato--Seiler--Simon inequality, variational formula for $\log\det_\zeta$) are standard.
\end{remark}
